\section{Algorithmic Typing}
\label{sec:algorithmic-typing}

This section develops an algorithmic version of the type system in
\cref{sec:types}. The algorithmic system is directly amenable to the
construction of a type checking algorithm. It employs a sound heuristic for
inferring the remainder session types of split operations. Type checking is a
two step process: first, a set of constraints is collected involving fresh
unification variables; second, the unification variables are solved and
validated with respect to the constraints. The constraint generation is part of
the algorithmic typing judgements in \cref{sec:alg-typing-relation}. Constraint
solving is described in \cref{sec:constraint-solving}.


\subsection{Algorithmic Judgments}
\label{sec:alg-typing-relation}

Algorithmic typing requires an extended grammar of session types
$\AS$ that includes unification variables $\TVar$ drawn from $\TVAR$. We write
$\AT$ for types that refer to $\AS$ instead of $\S$. Typing judgements collect
sets of constraints $\CSet$. Constraints $\CS$ come as either equality
constraints between two types, or mobility constraints:\quad$ \CS \grmeq* \AT
\CEq \ATU \grmor\CMobile \AT $

The premises of the algorithmic rules use a constraint-generating variant
$\Ctx[1] \SubCtx \Ctx[2] \uparrow \CSet$ of the subcontext relation. It is
defined by the rules of \cref{fig:syntax-types}, except that the two axioms
that transmute between parallel and sequential composition emit constraints
$\{\CMobile\AT \mid \CtxVar\EVar\AT \in \Ctx[i]\}$ instead of posing
$\Mobile{\Ctx[i]}$ as a requirement.
The variant is needed because algorithmic contexts may contain unification
variables: \RuleName{A-LetPair} and \RuleName{A-Case} bind variables at
inferred types. For example, the residual endpoint of a remote split has type
$\SSeq\SAcq\TVar$, which is mobile if the solution for $\TVar$ is
bounded. That is unknown when subcontext relation is checked. For every
$\sigma \models \CSet$, $\Ctx[1] \SubCtx \Ctx[2] \uparrow \CSet$ implies
$\Ctx[1][\sigma] \SubCtx \Ctx[2][\sigma]$. For contexts without unification
variables, the constraints can be solved immediately.

\begin{figure}[tp]
  \small
  \hfill \fbox{$\InferType \E \T \Eff$}
  \begin{mathpar}
    \RuleAlgVar \and
    \RuleAlgConst \and
    \RuleAlgAnnot \and
    \RuleAlgLSplit \and
    \RuleAlgRSplit \and
    \RuleAlgApp
    % \and \RuleAlgSeq
    \and \RuleAlgLetPair
    % \and \RuleAlgCase
  \end{mathpar}

  \medskip
  \setlength{\abovedisplayskip}{0pt}
  \setlength{\abovedisplayshortskip}{0pt}
  \setlength{\belowdisplayskip}{0pt}
  \setlength{\belowdisplayshortskip}{0pt}
  \hfill \fbox{$\CheckType \E \T \Eff$}
  \begin{mathpar}
    \RuleAlgAbs \and
    \RuleAlgCheck \and
    \RuleAlgAbsRec \and
    \RuleAlgPair
    %\and\RuleAlgInj 
  \end{mathpar}

  % Ensures equal space between judgements and definitions below as between
  % definitions and figure caption.
  \vspace{\abovecaptionskip}
  \[ \DefinitionsReqDirEff \]

  \caption{Selection of algorithmic typing rules}
  \label{fig:typing-algorithmic}
\end{figure}

\Cref{fig:typing-algorithmic} presents a selection of the bidirectional
algorithmic typing judgments~\cite{DBLP:journals/csur/DunfieldK21}. The
synthesis judgment $\InferType \E \AT \Eff$ takes as input the context $\Ctx$
and expression $\E$. It computes the expression's type $\AT$, its effect
$\Eff$, and a set of constraints $\CSet$. The checking judgment $\CheckType \E
\AT \Eff$ takes as input the context $\Ctx$, expression $\E$, and type $\AT$.
It returns the expression's effect $\Eff$ as well as a set of constraints
$\CSet$. Judgements make use of context join
$\Ctx[1] \CtxJoin*[\raisebox{-.4ex}{\kern1pt\relax$\scriptstyle\Dir$}] \Ctx[2]$, which
composes two contexts according to the direction $\Dir$, and context restriction
$\CtxRestrict*{X}$, which restricts a context to a set of variables $X$; the
definitions are given in \cref{fig:context-operators}.

In the synthesis rules, for variables the context must subsume to a single
variable binding (\RuleName{A-Var}) and to an empty context for constants
(\RuleName{A-Const}). A type annotation $\E:\T$ (\RuleName{A-Annot}) checks
$\E$ against $\T$ and synthesizes $\T$.
%
Function application \RuleName{A-App} infers the type of the function, checks
the argument against the expected type, verifies that the restricted contexts
join to a subcontext of $\Ctx$, and ensures the compatibility between the
involved directions and effects. Leftward application is compatibility only
with left-ordered functions; rightward application is compatible with all other
functions. The function's direction $\Dir*$ also determines which subterm may
have effects. If the function is $\DirLeft$-orderd the function term itself
must be pure; in case of a $\DirRight$-ordered function the argument term is
restricted to pure expressions.

The splitting constants receive special treatment by the algorithmic system.
Both constants have to be annotated with the borrowed type $\S$, which must not
be $\SSkip$ for $\CSplit$. \RuleName{A-LSplit} and \RuleName{A-RSplit} express
their synthesized types in terms of a fresh variable $\TVar$. The $\neg(\S[2]
\SEq \SSkip)$ and $\SLocal \S[2]$ requirements from the constant typing of
$\CSplit$ and $\CSSplit$ (see~\cref{fig:constant-typing}) need not be explicit
here. Only local, non-skip types are valid solutions to $\TVar$. The constants
themselves do not place any any constraints upon $\TVar$. These are generated
due to the interplay of \RuleName{A-App} and \RuleName{A-Check}: When
$\CSplit_{\S}$ is applied to an argument, the type of the argument is checked
against $\SSeq\S\TVar$.

Pair destructuring (\RuleName{A-LetPair}) in the algorithmic system must be
annotated with $\Dir\in\{\DirUnord,\DirLeft\}$. The such specified direction
concerns the ordering (or non-ordering) between the contexts of subexpressions
$\E[1]$ and $\E[2]$. The inferred type of $\E[1]$ comes with another direction
$\Dir*$ that indicates the ordering between the newly introduced bindings
$\EVar$ and $\EVary$. Sum type elimination works with a similar annotation.
\RuleName{A-Case} and the rule for sequencing (\RuleName{A-Seq}) are included
in the full set of inference judgements in
\cref{fig:appendix-algorithmic-infer}.

Rule \RuleName{A-LetPair} mirrors
its declarative counterparts: the context of the bound expression and the
context of the body are joined either in parallel or in sequence ($\CtxPatOp$),
and the new bindings enter the body's context with the same connective. The two
bindings are joined according to the order $\Dir$ of the pair. Elimination of a
sum type using \RuleName{A-Case} works similarly. The branches are typed in the
context restricted to the variables that they use. In addition, we have to add
equations to $\CSet$ that force all branch output types to be equal.

The checking rules cover lambda abstraction, pair construction, and sum
injection. For all other constructs \RuleName{A-Check} switches to inference
mode and emits an equality constraint between the given type $\T$ and the
inferred type $\TU$. For lambdas, \RuleName{A-Abs} enforces the supplied
function type, which specifies the mobility $\Mob$, the direction $\Dir$, and
the effect $\Eff$. If $\Mob = \MobMobile$, then all free variables in the
lambda's body must be mobile. The direction enforces the way the new binding is
joined to the context. The effect puts an upper bound on the effect of the
body. The mobility constraints are inserted into the outgoing constraint set
$\CSet$. \RuleName{A-AbsRec} is similar, but specialized to a recursive
function, which must be unrestricted and hence mobile. \RuleName{A-Pair} checks
the components pointwise with $\Ctx$ restricted to the set of variables
occuring free in the particular subexpression. An additional check verifies
that joining the partial contexts according to $\Dir$ gives a subcontext of
$\Ctx$. The premise $\reqEff{\Dir}{\Eff[2]}{\Eff[1]}$ ensures that if the pair
is ordered, depending on the order, one particular subterm is pure. Checking of
sum injections is standard and included in the full set checking judgements
(\cref{fig:appendix-algorithmic-check}).


\subsection{Constraint Solving}
\label{sec:constraint-solving}

An algorithmic typing for an expression $\E$ yields a set of
constraints $\CSet$. In inference mode, the returned type $\AT$ may refer to
unsolved unification variables. The goal of constraint solving is to find a
closing substitution. The typing is invalid if no such substitution exists.

\begin{definition}
  A substitution $\sigma : \TVAR \to \AS$ is a \emph{closing
  substitution} for  a set of constraints $\CSet$, written $\sigma \models
  \CSet$, if for every variable $\TVar$ in $\CSet$, $\sigma (\TVar) \in \S$,
  $\neg(\sigma(\TVar) \SEq \SSkip)$, and $\SLocal{\sigma(\TVar)}$, and all
  constraints are respected.
\end{definition}

Finding such a substitution boils down to a unification problem, but we are not
aware of a unification algorithm for the theory of context-free session types,
which involves associativity, neutral elements, and distributivity. Hence, we
approach constraint solving by constraint simplification followed by a sound
heuristic.

The equivalence constraints in $\CSet$ are simplified by structural
decomposition: pairs, sums, function types, and contexts are reduced
componentwise, while unit constraints disappear. The full rule set is given
in \cref{fig:constraint-simplification}. The rules are applied
exhaustively. Constraints of the form $\TVar\CEq\AT$ provide a potential
solution for unification variable $\TVar$. If left with unification variables
occurring in $\CSet$ or $\AT$ for which no solution candidate is found, the
heuristic fails and $\E$ is rejected.

Potential solutions have to first be identified and then validated with respect to the set $\CSet$.
We non-deterministically pick one candidate instantiation per variable, resulting in
a variable-free $\CSet*$\!. 
If there are multiple candidates for a unification variable the type to be substituted
into $\CSet$ can be picked non-deterministically. To see this, consider the case where after
simplification there are two constraints $\TVar \CEq \T$ and $\TVar \CEq \TU$. If
$\T$ is not a valid solution for $\TVar$, then $\TU$ is also not a solution.
The equality $\T \simeq \TU$ must hold for all types that can be selected. Since mobility
respects type-equivalence no mobility constraints exist validated by one but not the other candidate.

Mobility constraints can be
validated after solving the unification constraints because $\Mobile\T$ is decidable. 
Equality constraints can be verified because the equivalence of CFSTs is
decidable~\cite{DBLP:conf/tacas/AlmeidaMV20}. 
As our type language is richer than ordinary CFSTs (e.\,g. the distinction
between left and right-ordered function types), we require
an equivalence respecting encoding from our types to CFSTs.
\Cref{sec:algorithmic-typing-appendix} describes the
translation and sketches a correctness proof for the encoding. 

% In the following we sketch why our heuristic is able to infer useful
% information. Fresh unification variables are only introduced through
% \RuleName{A-LSplit} and \RuleName{A-RSplit}. In the case for $\CSplit$ the
% fresh variable $\TVar$ occurs nested in the function result type not composed
% with other session types. The function itself is linear, guaranteeing that a
% value of the result type will be lying around. \js{continue here}
% 
% $\SSkip$ as well as any other session types equivalent to skip \js{do we use this
% terminology?} are never valid solutions for unification variables, because this would
% be allow to split a session type off from itself.
% 
% One instance where the heuristic falls short is when the second component of a
% $\CSSplit$-result is split further with a prefix of more than just $\SAcq$.

\subsection{Implementation}

We supplement our theoretical work on the bidirectional type checking algorithm
from this section with an implementation written in Rust. It is based on an
artifact of prior work~\cite{saffrich25:_artif_borrow_from_session_types} and
uses
FreeST~5.0~\cite{DBLP:journals/corr/abs-1904-01284,DBLP:journals/iandc/AlmeidaMTV22,
almeida26:_frees_concur_progr_languag_with} for equivalence checking of session
types. The FreeST integration is described more closely in
\cref{sec:appendix-implementation}.

%The test cases of the Rust implementation include the motivational examples
%from \cref{sec:motivation}.

Because of the incompleteness of the heuristic from
\cref{sec:constraint-solving} the type checker may reject typeable programs.
The user can guide the constraint solver through additional type annotations.
Alongside six negative example, we provide 16~positive ones (including those
from \cref{sec:motivation}) where the heuristic works satisfactorily.


%%% Local Variables:
%%% mode: latex
%%% TeX-master: "../popl2027.tex"
%%% End:
